\begin{abstract}
Satellite Internet of Things services reach remote rivers but sell only a few short, one-way messages per day, and the server never learns which messages were lost. This study asked what a server can decide about regulatory thresholds from the bytes that arrive. A node sends a self-contained error-bounded record within the accuracy that the U.S. Geological Survey requires of an excellent record, plus witnesses: quantised daily extrema and a conductance rate pair. A three-valued certificate rule then answers exceeded, not exceeded, or undetermined, and the answers are sound under stated assumptions. The method was developed on 38 rivers (2023) and confirmed once on 31 sealed rivers (2024) against nine frozen hypotheses. Eight held. No certified decision was wrong, and 30,285 message blocks passed a byte-level round trip without a mismatch. Targeted witnesses raised decisions on hard cases over a record alone by 0.30 to 0.32 at six messages per day with 30\% loss. Decision flags beat witnesses on thresholds fixed at configuration time, but a hybrid that hosts witnesses in spare slots of the flag design cost nothing and added 0.29 to 0.36 on thresholds that nobody configured. The ninth hypothesis, a bound on false certificates after a forged message that passes the 16-bit check, failed at twelve messages per day. The evidence is trace-driven; no satellite hardware was used.
\end{abstract}

\begin{keyword}
water-quality monitoring \sep satellite Internet of Things \sep error-bounded records \sep certified decisions \sep erasure coding \sep pre-registered evaluation
\end{keyword}
